\documentclass{article}
\usepackage{graphicx} % Required for inserting images
\usepackage{amsmath}
\usepackage{amsfonts}                                                                   
\usepackage{amsthm}
\usepackage{amssymb}
\usepackage[utf8]{inputenc}
\usepackage[left=2cm, top=-5cm, right=2cm, bottom=2cm, includeheadfoot, headheight=130pt, a4paper]{geometry}
\usepackage[polish]{babel}
\usepackage[utf8]{inputenc}
\usepackage{polski}
\usepackage[T1]{fontenc}
\usepackage{longtable}
\usepackage{caption}
\usepackage{subcaption}
\usepackage{booktabs}
\usepackage{array}
\usepackage{float}
\usepackage{hyperref}
\usepackage{float}
\usepackage{longtable}

\usepackage[polish]{babel}
\usepackage[utf8]{inputenc}
\usepackage{polski}
\usepackage[T1]{fontenc}

\title{Algorytmy Metaheurystyczne\\Lista 0 - TSP}
\author{Rafał Włodarczyk}
\date{10.03.2026}

\begin{document}

\maketitle

\tableofcontents

\section{Źródło dla danych}

Dane do programu są wzięte ze strony internetowej:
\begin{center}
\texttt{https://www.math.uwaterloo.ca/tsp/world/countries.html}
\end{center}
oraz są dostępne do pobrania w formacie \texttt{.tsp} za pomocą skryptu \texttt{get\_data.sh}.

\section{Zadanie 1 - Metoda Losowa}

Zdefiniujmy odległość Euklidesową jako:
\begin{align*}
    d((x_1, y_1), (x_2, y_2)) = \sqrt{(x_2 - x_1)^2 + (y_2 - y_1)^2}
\end{align*}
W zadaniu użyłem \texttt{std::shuffle} do losowania permutacji wierzchołków.

\noindent
W pliku \texttt{task1.cpp} zapisany został kod źródłowy programu 
wczytującego z dane grafów oraz wyznaczającego statystyki 1000 permutacji,
które prezentują się następująco:

\begin{longtable}{|l|rrr|}
\hline
\textbf{Graf} & \textbf{Min} & \textbf{Min10Avg} & \textbf{Min50Avg} \\ \hline
western\_sahara.tsp & 77764 & 91586 & 82601 \\ 
djibouti.tsp & 20884 & 23528 & 21335 \\ 
qatar.tsp & 81614 & 86907 & 81237 \\ 
uruguay.tsp & 1555856 & 1581006 & 1495132 \\
zimbabwe.tsp & 2243300 & 2283067 & 2160078 \\ \hline
\end{longtable}
\noindent
Wartości zostały wyznaczane w precyzji \texttt{double} zaokrąglone do najbliższej liczby całkowitej.
Wyjaśnienie kolumn (wynik w jednostce wagi krawędzi - odległości w płaszczyźnie euklidesowej):
\begin{itemize}
    \item \textbf{Min} oznacza najmniejszą całkowitą wagę spośród 1000 permutacji
    \item \textbf{Min10Avg} średnia z liczb - najmniejsza całkowita waga w kolejnych 10 losowaniach.
    \item \textbf{Min50Avg} średnia z liczb - najmniejsza całkowita waga w kolejnych 50 losowaniach. 
\end{itemize}

\section{Zadanie 2 - Wykresy}

Rozwiązanie graficzne obejmuje stworzenie wykresu w płaszczyźnie Euklidesowej, w której
poszczególne punkty reprezentują wierzchołki grafu, a linie łączące te punkty reprezentują krawędzie,
łączymy krawędzie w kolejności permutacji uzyskanej w zadaniu 1, tworząc w ten sposób trasę.\\\\

\noindent
Przedstawiam wykresy dla każdego z grafów, które zostały wygenerowane w pliku \texttt{task2.cpp}:

\begin{figure}[H]
    \centering
    \begin{minipage}{0.49\textwidth}
        \centering
        \includegraphics[width=\textwidth]{output/best_randomized_western_sahara.tsp.png}
        \caption{Trasa dla grafu western\_sahara.tsp}
        \label{fig:western_sahara}
    \end{minipage}
    \hfill
    \begin{minipage}{0.49\textwidth}
        \centering
        \includegraphics[width=\textwidth]{output/best_randomized_djibouti.tsp.png}
        \caption{Trasa dla grafu djibouti.tsp}
        \label{fig:djibouti}
    \end{minipage}
    \vspace{1cm}
    \begin{minipage}{0.49\textwidth}
        \centering
        \includegraphics[width=\textwidth]{output/best_randomized_qatar.tsp.png}
        \caption{Trasa dla grafu qatar.tsp}
        \label{fig:qatar}
    \end{minipage}
    \hfill
    \begin{minipage}{0.49\textwidth}
        \centering
        \includegraphics[width=\textwidth]{output/best_randomized_uruguay.tsp.png}
        \caption{Trasa dla grafu uruguay.tsp}
        \label{fig:uruguay}
    \end{minipage}
    \vspace{1cm}
    \begin{minipage}{0.49\textwidth}
        \centering
        \includegraphics[width=\textwidth]{output/best_randomized_zimbabwe.tsp.png}
        \caption{Trasa dla grafu zimbabwe.tsp}
        \label{fig:zimbabwe}
    \end{minipage}
\end{figure}

\subsection{Wniosek}

\section{Zadanie 3 - MST}

Do wykonania zadanie posłużyłem się algorytmem Kruskala, który był szerzej omawiany na przedmiocie
Algorytmy i Struktury Danych. W pliku \texttt{task3.cpp} znajduje się kod źródłowy programu, 
który wczytuje dane grafów, wyznacza minimalne drzewa rozpinające (MST) oraz oblicza wagi tych drzew.\\\\

\noindent
Wyniki prezentują się następująco:
\begin{longtable}{|l|r|r|}
\hline
\textbf{Graf} & \textbf{Liczba wierzchołków (n)} & \textbf{Waga MST} \\ \hline
western\_sahara.tsp & 29 & 21751 \\ 
djibouti.tsp & 38 & 5831 \\ 
qatar.tsp & 194 & 8028 \\ 
uruguay.tsp & 734 & 69876 \\ 
zimbabwe.tsp & 929 & 82873 \\ \hline
\end{longtable}

\section{Zadanie 4 - MST do TSP z Wykresami}

Możemy wykonać DFS na MST, aby uzyskać trasę rozwiązującą TSP,
uzyskując w ten sposób 2-aproksymację rozwiązania. W pliku \texttt{task4.cpp} znajduje się kod źródłowy programu.\\\\

\noindent
Cykl komiwojażera powstaje w wyniku wypisania wierzchołków w kolejności odwiedzania podczas przeszukiwania w głąb.
To znaczy że wykorzystujemy MST do stworzenia grafu, po którym następnie wykonujemy DFS, a kolejność odwiedzania wierzchołków podczas DFS tworzy trasę TSP.\\\\

\begin{figure}[H]
    \centering
    \begin{minipage}{0.49\textwidth}
        \centering
        \includegraphics[width=\textwidth]{output/kruskal_western_sahara.tsp.png}
        \caption{Trasa dla grafu western\_sahara.tsp}
        \label{fig:mst-western_sahara}
    \end{minipage}
    \hfill
    \begin{minipage}{0.49\textwidth}
        \centering
        \includegraphics[width=\textwidth]{output/kruskal_djibouti.tsp.png}
        \caption{Trasa dla grafu djibouti.tsp}
        \label{fig:mst-djibouti}
    \end{minipage}
    \vspace{1cm}
    \begin{minipage}{0.49\textwidth}
        \centering
        \includegraphics[width=\textwidth]{output/kruskal_qatar.tsp.png}
        \caption{Trasa dla grafu qatar.tsp}
        \label{fig:mst-qatar}
    \end{minipage}
    \hfill
    \begin{minipage}{0.49\textwidth}
        \centering
        \includegraphics[width=\textwidth]{output/kruskal_uruguay.tsp.png}
        \caption{Trasa dla grafu uruguay.tsp}
        \label{fig:mst-uruguay}
    \end{minipage}
    \vspace{1cm}
    \begin{minipage}{0.49\textwidth}
        \centering
        \includegraphics[width=\textwidth]{output/kruskal_zimbabwe.tsp.png}
        \caption{Trasa dla grafu zimbabwe.tsp}
        \label{fig:mst-zimbabwe}
    \end{minipage}
\end{figure}

\section{Porównanie wyników}

Porównanie następuje pomiędzy metodą \textbf{Random} (zadanie 1), \textbf{MST 2-approx} (zadanie 4), \textbf{OPT} - informacje na stronie internetowej.

\begin{longtable}{|l|r|r|r|}
\hline
\textbf{Graf} & \textbf{Random} & \textbf{MST 2-approx} & \textbf{OPT} \\ \hline
western\_sahara.tsp & 77764 & 35009 & 27603 \\ 
djibouti.tsp & 20884 & 8739 & 6656 \\ 
qatar.tsp & 81614 & 12510 & 9352 \\ 
uruguay.tsp & 1555856 & 107340 & 79114 \\ 
zimbabwe.tsp & 2243300 & 131828 & 95345 \\ \hline
\end{longtable}
\noindent

Porównanie wyników pokazuje, że metoda MST 2-aproksymacji daje wyniki znacznie bliższe wartościom optymalnym (OPT) w porównaniu do metody losowej.
Widzimy również doświadczalnie, że dla każdego przypadku metody \textbf{MST 2-approx} wynik nie przekracza dwu-krotności wartości optymalnej,
co jest zgodne z teoretyczną gwarancją tej aproksymacji.

\end{document}